\chapter{Linear Algebra and Polynomials}\label{ch:linalg}

Two kinds of algebra appear again and again in contest problems. Linear algebra turns walks, linear systems and XOR puzzles into
matrix operations. Polynomials turn counting into multiplication: the coefficients of a product are a convolution, the fast Fourier
transform computes it in $\Oh{n \log n}$, and Newton's method lifts that to inverses, logarithms and exponentials of power series.
This chapter covers both, from matrix powers to the polynomial toolkit, and ends with transforms for bitwise convolutions and with
recovering linear recurrences from their first terms.

\begin{chapterintro}{\gold\ \plat\ \ioi}{Section~\ref{sec:matexp} (matrix powers for recurrences), Chapter~\ref{ch:numbertheory}
(modular inverses, primitive roots) and Section~\ref{sec:sosdp} (subset sums).}
  \item count and optimise walks with matrix powers over different semirings;
  \item solve linear systems modulo a prime with Gaussian elimination, and reason about XOR with bases over GF(2);
  \item multiply polynomials with the FFT and the NTT, and recognise convolutions in disguise;
  \item invert, divide, and take logarithms and exponentials of power series with Newton's method, and evaluate and interpolate at many
        points;
  \item compute XOR, AND and OR convolutions with the Walsh--Hadamard transform;
  \item find a linear recurrence from its first terms (Berlekamp--Massey) and jump to its $n$-th term (Kitamasa).
\end{chapterintro}

% ======================================================================================================================
\section[Matrices and walks]{Matrices and walks \hfill \gold}\label{sec:matrices}

Multiplying an $n \times m$ matrix by an $m \times p$ matrix costs $nmp$ multiplications, so a power $A^k$ of an $n \times n$ matrix
costs $\Oh{n^3 \log k}$ by repeated squaring. Section~\ref{sec:matexp} used this for linear recurrences. The other classic use is
counting walks: if $A$ is the adjacency matrix of a graph, then $(A^k)_{uv}$ is the number of walks with exactly $k$ edges from $u$ to
$v$, because a walk of length $k$ is a walk of length $k - 1$ followed by one more edge, which is exactly how a matrix product combines
its entries (Figure~\ref{fig:walks}).

\begin{figure}[htbp]
\centering
\begin{tikzpicture}
  \tikzset{wv/.style={vtx, minimum size=7mm}}
  \node[wv] (W1) at (0, 1.2) {1};
  \node[wv] (W2) at (1.6, 1.2) {2};
  \node[wv] (W3) at (0.8, 0) {3};
  \draw[ded, shorten >=1pt] (W1) -- (W2);
  \draw[ded, shorten >=1pt] (W2) -- (W3);
  \draw[ded, shorten >=1pt] (W1) to[bend right=18] (W3);
  \draw[ded, shorten >=1pt] (W3) to[bend right=18] (W1);
  \node[anchor=west, font=\small] at (2.4, 0.6) {$A = \begin{pmatrix} 0 & 1 & 1 \\ 0 & 0 & 1 \\ 1 & 0 & 0 \end{pmatrix}, \quad
    A^2 = \begin{pmatrix} 1 & 0 & 1 \\ 1 & 0 & 0 \\ 0 & 1 & 1 \end{pmatrix}$};
  \node[notel, anchor=west, text width=4.6cm] at (8.6, 0.6) {Row 1 of $A^2$: from vertex 1 there is one walk of two edges back to
    1 (via 3) and one to 3 (via 2), and none to 2.};
\end{tikzpicture}
\caption{The entries of $A^k$ count the walks with $k$ edges.}
\label{fig:walks}
\end{figure}

\usetemplate{Matrix exponentiation}{Ctrl+Alt+T, Shift+X}{cpmatrix}

\begin{problem}{Graph Paths I (CSES)}
Count the routes with exactly $k \le 10^9$ flights from city 1 to city $n$ ($n \le 100$, parallel flights allowed), modulo $10^9 + 7$.
\solution $(A^k)_{1n}$ with $A_{uv}$ the number of flights from $u$ to $v$: 30 squarings of a $100 \times 100$ matrix, 0.11\,s.
\end{problem}

\begin{code}
void solve() {  // Mat, mat_pow (cpmatrix) are pasted above solve()
    int n, m, k;
    cin >> n >> m >> k;
    Mat A(n, vector<ll>(n));
    rep(i, m) {
        int a, b;
        cin >> a >> b;
        A[a - 1][b - 1]++;  // parallel flights count separately
    }
    print(mat_pow(A, k)[0][n - 1]);  // (A^k)[u][v]: walks with exactly k edges
}
\end{code}

\paragraph{Other semirings.} Nothing in repeated squaring needs ordinary $+$ and $\times$: any pair of operations in which $\otimes$
distributes over $\oplus$ and both are associative works. With $(\min, +)$, the ``product'' $\min_l (X_{il} + Y_{lj})$ combines the
cheapest ways to go from $i$ to $l$ and from $l$ to $j$, and $A^k$ holds the cheapest walks with exactly $k$ edges; its identity matrix
has 0 on the diagonal and $\infty$ elsewhere. With $(\max, \min)$ it gives the best bottleneck, and with $(\lor, \land)$ plain
reachability in exactly $k$ steps, where rows stored as bitsets make each product 64 times faster.

\begin{problem}{Graph Paths II (CSES)}
The cheapest route with exactly $k \le 10^9$ flights from city 1 to city $n$ ($n \le 100$), or $-1$.
\solution The $(\min, +)$ power of the price matrix. Prices up to $10^9$ times $10^9$ flights stay below $10^{18}$, so \cc{INF}
$= 4 \cdot 10^{18}$ is safe. 0.09\,s.
\end{problem}

\begin{code}
void solve() {
    int n, m, k;
    cin >> n >> m >> k;
    vvi A(n, vi(n, INF));  // A[u][v]: the cheapest single flight from u to v
    rep(i, m) {
        int a, b;
        ll c;
        cin >> a >> b >> c;
        chmin(A[a - 1][b - 1], c);
    }
    auto mul = [&](const vvi& X, const vvi& Y) {  // the (min, +) product
        vvi Z(n, vi(n, INF));
        rep(i, n) rep(l, n) {
            if (X[i][l] == INF) continue;
            rep(j, n) if (Y[l][j] < INF) chmin(Z[i][j], X[i][l] + Y[l][j]);
        }
        return Z;
    };
    vvi R(n, vi(n, INF));  // the identity for (min, +): 0 on the diagonal
    rep(i, n) R[i][i] = 0;
    for (; k > 0; k >>= 1, A = mul(A, A))
        if (k & 1) R = mul(R, A);
    print(R[0][n - 1] < INF ? R[0][n - 1] : -1);
}
\end{code}

% ======================================================================================================================
\section[Gaussian elimination]{Gaussian elimination \hfill \plat}\label{sec:gauss}

To solve $n$ linear equations in $m$ unknowns, write them as an augmented matrix $(A \mid b)$ and simplify it with three row operations
that do not change the solutions: swapping rows, scaling a row by a non-zero number, and adding a multiple of one row to another. Going
column by column, choose a \emph{pivot} row with a non-zero entry in the column, scale it so the pivot becomes 1, and subtract multiples
of it from every other row to clear the rest of the column (Figure~\ref{fig:gauss}). This \emph{Gauss--Jordan} form shows the answer
directly; each column costs $\Oh{nm}$, so the whole elimination is $\Oh{nm \min(n, m)}$.

\begin{figure}[htbp]
\centering
\newcommand{\gpiv}[1]{\textcolor{fgverm}{\mathbf{#1}}}
$\left(\begin{array}{rrr|r} \gpiv{1} & 2 & 1 & 8 \\ 2 & 5 & 1 & 15 \\ 1 & 3 & 3 & 16 \end{array}\right)
\xrightarrow{\text{column 1}}
\left(\begin{array}{rrr|r} 1 & 2 & 1 & 8 \\ 0 & \gpiv{1} & -1 & -1 \\ 0 & 1 & 2 & 8 \end{array}\right)
\xrightarrow{\text{column 2}}
\left(\begin{array}{rrr|r} 1 & 0 & 3 & 10 \\ 0 & 1 & -1 & -1 \\ 0 & 0 & \gpiv{3} & 9 \end{array}\right)
\xrightarrow{\text{column 3}}
\left(\begin{array}{rrr|r} 1 & 0 & 0 & 1 \\ 0 & 1 & 0 & 2 \\ 0 & 0 & 1 & 3 \end{array}\right)$
\caption{Gauss--Jordan elimination on $x + 2y + z = 8$, $2x + 5y + z = 15$, $x + 3y + 3z = 16$. Each step uses the pivot (red) to clear
its column; the last matrix reads $x = 1$, $y = 2$, $z = 3$.}
\label{fig:gauss}
\end{figure}

The final form also tells what happens in general. The number of pivots is the \defn{rank} $r$. A row that has become $0 = c$ with
$c \ne 0$ means the system has no solution. Otherwise, the $m - r$ columns without a pivot are \emph{free variables}: any values work,
and the pivot variables follow. Modulo a prime $p$ there are then exactly $p^{m-r}$ solutions; the template returns one of them, with the
free variables set to 0. The same elimination computes determinants (the product of the pivots, with a sign for every swap,
Section~\ref{sec:matrixtree}) and inverse matrices (eliminate $(A \mid I)$ to $(I \mid A^{-1})$).

\usetemplate{Gaussian elimination mod $p$}{Ctrl+Alt+U, N}{cpgauss}

\begin{problem}{System of Linear Equations (CSES)}
Solve $n \le 500$ linear equations in $m \le 500$ unknowns modulo $10^9 + 7$, printing any solution, or $-1$ if there is none.
\solution \cc{gauss_mod}; scaling by a pivot uses its modular inverse. $500 \times 500$ takes 0.15\,s.
\end{problem}

\begin{code}
void solve() {  // gauss_mod (cpgauss) is pasted above solve()
    int n, m;
    cin >> n >> m;
    vvi a(n, vi(m));
    vi b(n), x;
    rep(i, n) {
        for (ll& v : a[i]) cin >> v;
        cin >> b[i];
    }
    if (gauss_mod(a, b, x) == -1) print(-1);  // inconsistent
    else print(x);  // free variables are 0
}
\end{code}

\begin{pitfall}[Elimination with floating point]
Over the reals, a tiny pivot amplifies rounding errors. Choose as pivot the entry of largest absolute value in the column (\emph{partial
pivoting}), compare with a tolerance such as $10^{-9}$ instead of 0, and prefer exact arithmetic (modulo a prime, or fractions) whenever
the problem allows it.
\end{pitfall}

\begin{recognize}
``Each switch toggles some lamps; can all lamps be lit?'' (a system over GF(2): rows are bitsets and elimination is 64 times faster with
XOR of whole words), expected values on a graph with cycles (one equation $E_v = 1 + \sum_w p_{vw} E_w$ per state), and anything phrased
as ``linear combinations'', ``independent vectors'' or ``a unique solution''.
\end{recognize}

% ======================================================================================================================
\section[Linear bases over GF(2)]{Linear bases over GF(2) \hfill \ioi}\label{sec:xorbasis}

Numbers under XOR form a vector space over the two-element field: a number is a vector of bits, and XOR adds vectors. Questions about
``the XOR of some subset'' are therefore questions about the \emph{span} of the numbers. A \defn{linear basis} keeps at most one vector
for each highest bit. To insert $x$, go through its bits from the top: whenever the basis has a vector with the same highest bit, XOR it
into $x$. If $x$ becomes 0, it was already representable; otherwise its remaining highest bit is new and $x$ joins the basis
(Figure~\ref{fig:xorbasis}).

\begin{figure}[htbp]
\centering
\begin{tikzpicture}
  \tikzset{bz/.style={cell, font=\ttfamily\small, minimum width=6mm, minimum height=6mm}}
  \foreach \lbr/\bits/\yy in {{$b_3$}/{1/cool,0/n,1/n,1/n}/0, {$b_2$}/{0/n,1/cool,1/n,0/n}/-0.75, {$b_0$}/{0/n,0/n,0/n,1/cool}/-1.5} {
    \node[blbl, anchor=east] at (-0.2, \yy) {\lbr};
    \foreach \bv/\bst [count=\ii from 0] in \bits \node[bz, \bst] at (0.6*\ii, \yy) {\bv};
  }
  \foreach \ii/\bt in {0/3, 1/2, 2/1, 3/0} \node[lbl] at (0.6*\ii + 0.3, 0.55) {\bt};
  \node[lbl, anchor=east] at (-0.2, 0.55) {bit};
  \node[notel, anchor=west, text width=8.2cm] at (3.1, -0.75) {Three vectors with distinct highest bits (blue) span $2^3 = 8$ values.
    The largest subset XOR is found greedily from the top: start from 0, take $b_3$ (1011), then $b_2$ because $1011 \oplus 0110 = 1101$
    is larger, and skip $b_0$ because $1101 \oplus 0001 = 1100$ is smaller. The maximum is $1101 = 13$.};
\end{tikzpicture}
\caption{A linear basis over GF(2): one vector per highest bit.}
\label{fig:xorbasis}
\end{figure}

With a basis of dimension $d$, the subset XORs take exactly $2^d$ distinct values, each reached by the same number $2^{n - d}$ of
subsets. The maximum subset XOR is the greedy of Figure~\ref{fig:xorbasis}, and ``can $x$ be formed?'' is an insertion that is not
stored. After a reduction that makes every basis vector's highest bit appear in no other vector, the $k$-th smallest subset XOR is
read off the binary representation of $k$.

\usetemplate{XOR basis}{Ctrl+Alt+Y, Shift+X}{cpxorbasis}

\begin{problem}{Maximum subset XOR}
For $n \le 2 \cdot 10^5$ numbers below $2^{60}$, print the largest XOR of a subset and the number of distinct subset XORs.
\solution Insert everything into a basis: $\Oh{n \cdot 60}$, 0.12\,s.
\end{problem}

\begin{code}
void solve() {  // XorBasis (cpxorbasis) is pasted above solve()
    int n;
    cin >> n;
    XorBasis xb;
    rep(i, n) {
        ll x;
        cin >> x;
        xb.insert(x);
    }
    print(xb.max_xor(), 1LL << xb.dim);  // 2^dim distinct subset XORs
}
\end{code}

\paragraph{Bases of ranges.} ``The largest subset XOR of $a_l, \dots, a_r$'' for many ranges, online, has a neat solution. Build a basis
of every prefix, but let each basis vector remember the position it came from, and when inserting a new element prefer it to an older
vector with the same highest bit: swap them and continue inserting the older one. Then, for the prefix ending at $r$, the basis vectors
whose positions are at least $l$ span exactly the XORs of subsets of $a_l, \dots, a_r$, and the greedy runs on them alone. Each prefix
costs one insertion, $\Oh{B}$ for $B$-bit values, and so does each query.

\begin{problem}{Range XOR maximum}
Answer $q \le 2 \cdot 10^5$ queries ``the largest subset XOR of $a_l, \dots, a_r$'' for $n \le 2 \cdot 10^5$ values below $2^{30}$.
\solution The prefix bases above, $30$ numbers and $30$ positions per prefix (48\,MB as 32-bit integers): 0.21\,s.
\end{problem}

\begin{code}
void solve() {
    int n, q;
    cin >> n >> q;
    const int B = 30;  // the values are below 2^30
    // prefix i: a basis of a[0..i) built from the latest possible positions
    vector<array<int32_t, B>> val(n + 1), pos(n + 1);
    val[0].fill(0), pos[0].fill(-1);
    rep(i, n) {
        int32_t x, p = i;
        cin >> x;
        val[i + 1] = val[i], pos[i + 1] = pos[i];
        per(b, B) {
            if (!(x >> b & 1)) continue;
            if (!val[i + 1][b]) {
                val[i + 1][b] = x, pos[i + 1][b] = p;
                break;
            }
            if (pos[i + 1][b] < p) {  // keep the newer vector in the basis
                swap(val[i + 1][b], x);
                swap(pos[i + 1][b], p);
            }
            x ^= val[i + 1][b];  // go on with the older of the two vectors
        }
    }
    while (q--) {
        int l, r;
        cin >> l >> r;
        int32_t best = 0;  // greedy from the top, vectors from a[l-1..r) only
        per(b, B) {
            int32_t v = val[r][b];
            if (pos[r][b] >= l - 1 && (best ^ v) > best) best ^= v;
        }
        print(best);
    }
}
\end{code}

\begin{recognize}
``Maximum XOR of a subset'', ``how many distinct XORs'', ``can $x$ be written as a XOR of some of these'', ``the $k$-th smallest subset
XOR'', and XOR along paths in a graph: the XOR of any cycle can be added to a path freely, so the cycles' XORs go into a basis and the
answer is the greedy maximum starting from one fixed path.
\end{recognize}

% ======================================================================================================================
\section[The fast Fourier transform]{The fast Fourier transform \hfill \ioi}\label{sec:fft}

The product of $A(x) = \sum a_i x^i$ and $B(x) = \sum b_j x^j$ has the coefficients $c_k = \sum_{i + j = k} a_i b_j$, the
\defn{convolution} of the two coefficient sequences (Figure~\ref{fig:convolution}). Computed directly it costs $\Oh{nm}$. Convolutions
are everywhere: whenever two independent choices add up, such as ``choose one item from each of two lists so that the weights sum to
$w$'', the counts of the sums are the convolution of the two histograms.

\begin{figure}[htbp]
\centering
\begin{tikzpicture}
  \tikzset{cv/.style={draw=fgline, line width=0.6pt, minimum size=9mm, inner sep=0pt, font=\sffamily\small},
    k0/.style={cv, fill=fgsky!35}, k1/.style={cv, fill=fgorange!30}, k2/.style={cv, fill=fggreen!28}, k3/.style={cv, fill=fgpink!35}}
  \foreach \av [count=\ii from 0] in {1, 2, 3} \node[lbl, anchor=east] at (-0.55, -0.95*\ii) {$a_{\ii} = \av$};
  \foreach \bv [count=\jj from 0] in {4, 5} \node[lbl] at (0.95*\jj, 0.75) {$b_{\jj} = \bv$};
  \foreach \ii/\jj/\pv in {0/0/4, 0/1/5, 1/0/8, 1/1/10, 2/0/12, 2/1/15} {
    \pgfmathtruncatemacro{\kk}{\ii + \jj}
    \node[k\kk] at (0.95*\jj, -0.95*\ii) {\pv};
  }
  \foreach \kk/\cv in {0/4, 1/13, 2/22, 3/15} \node[k\kk, minimum width=12mm] at (3.2 + 1.35*\kk, -0.95) {$c_{\kk} = \cv$};
  \node[notel, anchor=west, text width=4.4cm] at (8.7, -0.95) {$(1 + 2x + 3x^2)(4 + 5x)$: the products $a_i b_j$ with the same $i +
    j$ (same colour) add up to $c_k$, giving $4 + 13x + 22x^2 + 15x^3$.};
\end{tikzpicture}
\caption{A convolution sums the products along anti-diagonals.}
\label{fig:convolution}
\end{figure}

The \defn{fast Fourier transform} computes a convolution in $\Oh{n \log n}$ by changing representation. A polynomial of degree less
than $N$ is determined by its values at $N$ points, and in the value representation multiplication is pointwise. So:
\begin{enumerate}
  \item \textbf{evaluate} $A$ and $B$ at $N$ points, where $N$ is a power of two at least $n + m - 1$;
  \item \textbf{multiply} the values pointwise, $C(x_k) = A(x_k) B(x_k)$;
  \item \textbf{interpolate} $C$ back from its $N$ values.
\end{enumerate}
The trick is the choice of points: the complex $N$-th roots of unity $\omega^k = e^{2\pi i k / N}$ (Figure~\ref{fig:roots}). Split $A$
into its even and odd coefficients, $A(x) = A_{\mathrm{even}}(x^2) + x \cdot A_{\mathrm{odd}}(x^2)$. The squares of the $N$-th roots of
unity are the $N/2$-th roots of unity, each taken twice, so evaluating $A$ at all $N$ points needs only the two half-size polynomials at
$N/2$ points, plus $\Oh{N}$ work to combine them: $A(\omega^k) = E(\omega^{2k}) + \omega^k O(\omega^{2k})$ and $A(\omega^{k + N/2}) =
E(\omega^{2k}) - \omega^k O(\omega^{2k})$. Hence $T(N) = 2T(N/2) + \Oh{N} = \Oh{N \log N}$. Interpolation is the same transform with
$\omega^{-1}$ in place of $\omega$, followed by a division by $N$.

\begin{figure}[htbp]
\centering
\begin{tikzpicture}[scale=1.25]
  \draw[line width=0.4pt, draw=fgline] (-1.45, 0) -- (1.45, 0) (0, -1.45) -- (0, 1.45);
  \draw[line width=0.9pt, draw=fgblue!70] (0, 0) circle (1);
  \foreach \kk in {0, ..., 7} {
    \pgfmathtruncatemacro{\ev}{mod(\kk, 2) == 0 ? 1 : 0}
    \ifnum\ev=1
      \fill[fgblue] ({cos(45*\kk)}, {sin(45*\kk)}) circle (2.2pt);
    \else
      \fill[fgorange] ({cos(45*\kk)}, {sin(45*\kk)}) circle (2.2pt);
    \fi
    \node[font=\sffamily\scriptsize] at ({1.27*cos(45*\kk)}, {1.27*sin(45*\kk)}) {$\omega^{\kk}$};
  }
  \node[notel, anchor=west, text width=6.2cm] at (1.9, 0) {The eight 8th roots of unity. Squaring doubles the angle, so $\omega^k$
    and $\omega^{k+4} = -\omega^k$ have the same square, and the squares are exactly the even powers (blue), the four 4th roots of
    unity. Halving the number of points at every level is what makes the recursion fast.};
\end{tikzpicture}
\caption{The evaluation points of the FFT.}
\label{fig:roots}
\end{figure}

Implementations unroll the recursion. Reordering the input by bit-reversed indices puts the coefficients of every sub-problem next to
each other; then $\log_2 N$ passes combine neighbouring blocks of sizes $1, 2, 4, \dots$ with \emph{butterflies} $(u, v) \mapsto (u +
\omega^j v, u - \omega^j v)$ (Figure~\ref{fig:butterfly}).

\begin{figure}[htbp]
\centering
\begin{tikzpicture}[x=1.9cm, y=0.58cm]
  \foreach \sg/\half/\col in {1/1/fgblue, 2/2/fggreen, 3/4/fgverm} {
    \foreach \ii in {0, ..., 7} {
      \pgfmathtruncatemacro{\pp}{mod(floor(\ii / \half), 2) == 0 ? \ii + \half : \ii - \half}
      \draw[line width=0.7pt, draw=\col] (\sg - 1, -\ii) -- (\sg, -\ii);
      \draw[line width=0.7pt, draw=\col] (\sg - 1, -\pp) -- (\sg, -\ii);
    }
    \node[font=\sffamily\scriptsize\bfseries, text=\col] at (\sg - 0.5, 0.9) {blocks of \the\numexpr 2*\half\relax};
  }
  \foreach \sg in {0, ..., 3} \foreach \ii in {0, ..., 7} \fill[fgink] (\sg, -\ii) circle (1.7pt);
  \foreach \cf [count=\ii from 0] in {0, 4, 2, 6, 1, 5, 3, 7} \node[lbl, anchor=east] at (-0.06, -\ii) {$a_{\cf}$};
  \foreach \ii in {0, ..., 7} \node[lbl, anchor=west] at (3.06, -\ii) {$A(\omega^{\ii})$};
  \node[notel, anchor=west, text width=4.4cm] at (3.75, -3.5) {An 8-point FFT. The coefficients enter in bit-reversed order (4 =
    $100_2$ sits at position $001_2$). Each pass joins pairs of blocks; a node combines the two values that reach it, one of them
    multiplied by a root of unity.};
\end{tikzpicture}
\caption{The butterfly network of the iterative FFT: three passes for $N = 8$.}
\label{fig:butterfly}
\end{figure}

\usetemplate{FFT (exact integer convolution) / NTT}{Ctrl+Alt+Y, F / Shift+T}{cpfft, cpntt}

The library's \cc{multiply} packs both inputs into one complex array, $a + i b$, and needs a single forward and a single inverse
transform: squaring $a + ib$ gives $a^2 - b^2 + 2iab$, and the imaginary part of the product is twice the convolution.

\begin{problem}{Apples and Bananas (CSES)}
$n$ apples and $m$ bananas have weights between 1 and $k$ ($k, n, m \le 2 \cdot 10^5$). For every total $w$ from 2 to $2k$, count the
(apple, banana) pairs that weigh $w$ together.
\solution The convolution of the two weight histograms. Counts reach $4 \cdot 10^{10}$, well inside the FFT's exact range. 0.18\,s.
\end{problem}

\begin{code}
void solve() {  // multiply (cpfft) is pasted above solve()
    int k, n, m;
    cin >> k >> n >> m;
    vi a(k + 1), b(k + 1);  // histograms of the weights
    rep(i, n) {
        int x;
        cin >> x;
        a[x]++;
    }
    rep(i, m) {
        int x;
        cin >> x;
        b[x]++;
    }
    vi c = multiply(a, b);  // c[w] = sum of a[i] * b[w - i]
    print(vi(c.begin() + 2, c.begin() + 2 * k + 1));
}
\end{code}

\paragraph{Correlations.} Reversing one sequence turns a convolution, which pairs indices with a fixed \emph{sum}, into a correlation,
which pairs indices with a fixed \emph{difference}: with $b_i = a_{n-1-i}$, the coefficient $c_{n-1+k}$ of $a * b$ is $\sum_i a_{i+k}
a_i$.

\begin{problem}{One Bit Positions (CSES)}
For a bit string of length $n \le 2 \cdot 10^5$ and every distance $k$ from 1 to $n - 1$, count the pairs of 1-bits at distance $k$.
\solution The correlation of the string with itself: one FFT, 0.13\,s.
\end{problem}

\begin{code}
void solve() {  // multiply (cpfft) is pasted above solve()
    string s;
    cin >> s;
    int n = sz(s);
    vi a(n), b(n);
    rep(i, n) a[i] = b[n - 1 - i] = s[i] - '0';  // b is a reversed
    vi c = multiply(a, b);  // c[n - 1 + k] = sum over i of a[i + k] * a[i]
    vi ans(n - 1);
    rep(k, 1, n) ans[k - 1] = c[n - 1 + k];
    print(ans);
}
\end{code}

\begin{table}[htbp]
\centering\small
\begin{tabularx}{\linewidth}{@{}>{\raggedright\arraybackslash}p{6.4cm}X@{}}
\toprule
problem & the convolution \\
\midrule
pairs $(i, j)$ with $a_i + b_j = s$, for every $s$ & the histograms of $a$ and $b$ \\
pairs at distance $k$, or all pairwise differences & a histogram with the reversed other histogram \\
string matching with wildcards & $\sum_j (p_j - t_{i+j})^2 p_j t_{i+j}$, expanded into three convolutions \\
the number of mismatches of a pattern at every shift & one correlation per letter of the alphabet \\
one item from each of $k$ groups, by total weight & the product of the groups' generating polynomials \\
multiplying huge integers & digits as coefficients, then propagate the carries \\
\bottomrule
\end{tabularx}
\caption{Convolutions in disguise.}
\label{tab:convapps}
\end{table}

\begin{pitfall}[Precision of the complex FFT]
A \cc{double} has 53 bits of precision. The library's FFT computes its roots in long double and stays exact while the true
coefficients of the result are below about $10^{15}$; beyond that, results are silently rounded. For larger values, or values modulo an
arbitrary number, split each coefficient into two halves of 15 bits (a few more transforms), or use the NTT with three primes and the
Chinese remainder theorem.
\end{pitfall}

% ======================================================================================================================
\section[The number theoretic transform]{The number theoretic transform \hfill \ioi}\label{sec:ntt}

The FFT only needs a primitive $N$-th root of unity, and the integers modulo a prime $p$ have one whenever $N$ divides $p - 1$: if $g$ is
a primitive root (Section~\ref{sec:dlog}), then $g^{(p-1)/N}$ has order exactly $N$. For $p = 998\,244\,353 = 119 \cdot 2^{23} + 1$
every power of two up to $2^{23}$ divides $p - 1$, and $g = 3$. The \defn{number theoretic transform} is the FFT over these integers:
exact, with no rounding at all, and producing the coefficients modulo $p$, which is precisely what counting problems ask for. That is why
so many counting problems use $998\,244\,353$ instead of $10^9 + 7$, which has no large power of two in $p - 1$.

\begin{trick}[Any modulus]
For a convolution modulo $10^9 + 7$ (or exact results up to about $10^{26}$), run the NTT modulo three primes of the form $c \cdot 2^k +
1$, such as $998\,244\,353$, $167\,772\,161$ and $469\,762\,049$, and recover every coefficient by the Chinese remainder theorem: the
true value is below the product of the three primes, about $7.8 \cdot 10^{25}$.
\end{trick}

% ======================================================================================================================
\section[Operations on power series]{Operations on power series \hfill \ioi}\label{sec:polyops}

Counting problems speak the language of \emph{generating functions}: a sequence $f_0, f_1, \dots$ becomes the formal power series $F(x) =
\sum f_n x^n$, and operations on sequences become algebra. Multiplying two series convolves them. The ways to make a total of $n$ from
coins of values $c_1, \dots, c_k$, each usable any number of times, are the coefficients of $\prod_i \frac{1}{1 - x^{c_i}}$. With a fast
product, the other operations follow from \textbf{Newton's method}, and each costs $\Oh{n \log n}$ for the first $n$ coefficients.

\paragraph{The inverse.} If $A(x) B(x) \equiv 1 \pmod{x^k}$, then $B' = B (2 - A B)$ satisfies $A B' \equiv 1 \pmod{x^{2k}}$: writing
$AB = 1 - E$ with $E \equiv 0 \pmod{x^k}$, we get $A B' = (1 - E)(1 + E) = 1 - E^2$, and $E^2 \equiv 0 \pmod{x^{2k}}$. Every step doubles
the number of correct coefficients (Figure~\ref{fig:newton}) and costs a constant number of products of size $2k$, so the total is
$\Oh{n \log n}$, dominated by the last step. The inverse exists exactly when $a_0 \ne 0$.

\begin{figure}[htbp]
\centering
\begin{tikzpicture}
  \tikzset{nc/.style={cell, font=\sffamily\small, minimum width=7mm}}
  \node[blbl, anchor=east] at (-0.2, 0) {mod $x$};
  \arrs{Na}{0}{0}{1/good}
  \node[blbl, anchor=east] at (-0.2, -0.8) {mod $x^2$};
  \arrs{Nb}{0}{-0.8}{1/n, 1/good}
  \node[blbl, anchor=east] at (-0.2, -1.6) {mod $x^4$};
  \arrs{Nc}{0}{-1.6}{1/n, 1/n, 2/good, 3/good}
  \node[blbl, anchor=east] at (-0.2, -2.4) {mod $x^8$};
  \arrs{Nd}{0}{-2.4}{1/n, 1/n, 2/n, 3/n, 5/good, 8/good, 13/good, 21/good}
  \node[notel, anchor=west, text width=7.0cm] at (6.1, -1.2) {Inverting $A = 1 - x - x^2$, whose inverse generates the Fibonacci
    numbers. Starting from $B = 1$, each step $B \gets B (2 - AB)$ keeps the correct coefficients and adds as many new ones (green):
    $1 + x$, then $1 + x + 2x^2 + 3x^3$, then eight coefficients.};
\end{tikzpicture}
\caption{Newton's method doubles the number of correct coefficients at every step.}
\label{fig:newton}
\end{figure}

\paragraph{The rest of the toolkit.} The same doubling builds everything else from the inverse:
\begin{itemize}
  \item \textbf{Division with remainder}, $A = BQ + R$ with $\deg R < \deg B$: reversing the coefficient order turns the quotient into a
        power series division, $\mathrm{rev}(Q) = \mathrm{rev}(A) / \mathrm{rev}(B) \bmod x^{\deg A - \deg B + 1}$; then $R = A - BQ$.
  \item \textbf{Logarithm} (for $a_0 = 1$): $\log A = \int A' / A$, a derivative, an inverse, a product and an integral.
  \item \textbf{Exponential} (for $a_0 = 0$): Newton's method on $\log B - A = 0$ gives the step $B \gets B (1 - \log B + A)$.
  \item \textbf{Powers and roots}: $A^k = \exp(k \log A)$ after factoring out the lowest term, and $\sqrt A$ by $B \gets (B + A/B) / 2$.
\end{itemize}

\usetemplate{Polynomial toolkit mod 998244353}{Ctrl+Alt+U, Shift+P}{cppoly}

\begin{problem}{Partition numbers}
Print the number of partitions $p(1), \dots, p(n)$ modulo $998\,244\,353$, for $n \le 2 \cdot 10^5$.
\solution $\sum p(n) x^n = \prod_{k \ge 1} \frac{1}{1 - x^k}$, and Euler's pentagonal number theorem gives the denominator explicitly:
$\prod_k (1 - x^k) = \sum_j (-1)^j x^{j(3j-1)/2}$, summed over all integers $j$, a series with only $\Oh{\sqrt n}$ non-zero terms. One
inverse finishes the job: 0.34\,s, against $\Oh{n \sqrt n}$ for the classic DP over the pentagonal recurrence.
\end{problem}

\begin{code}
void solve() {  // poly_inv (cppoly) is pasted above solve()
    int n;
    cin >> n;
    Poly e(n + 1, 0);  // Euler: prod (1 - x^k) is a sparse series with +-1
    for (ll j = 0; j * (3 * j - 1) / 2 <= n; j++)  // at the pentagonal numbers
        for (ll g : {j * (3 * j - 1) / 2, j * (3 * j + 1) / 2})
            if (g <= n) e[g] = j % 2 ? PMOD - 1 : 1;
    Poly p = poly_inv(e, n + 1);  // the partition numbers: 1 / prod (1 - x^k)
    print(vi(p.begin() + 1, p.end()));
}
\end{code}

\paragraph{Products of many factors.} A product of many sparse factors $\prod (1 + x^{w_i})$ is best computed as the exponential of a sum
of logarithms, because $\log(1 + x^w) = x^w - \frac{x^{2w}}{2} + \frac{x^{3w}}{3} - \cdots$ is known in closed form. Grouping equal
weights, the sum of the logarithms costs $\sum_w S / w = \Oh{S \log S}$ terms by the harmonic bound, and one exponential follows.

\begin{problem}{Subset sums}
For $n \le 2 \cdot 10^5$ items with weights between 1 and $S \le 10^5$, count, for every $s$ from 1 to $S$, the subsets of total weight
exactly $s$, modulo $998\,244\,353$.
\solution $\prod_i (1 + x^{w_i}) = \exp\bigl(\sum_w \mathit{cnt}_w \log(1 + x^w)\bigr)$. The knapsack DP would need $\Oh{nS}$, $2 \cdot
10^{10}$ steps; this takes 0.41\,s.
\end{problem}

\begin{code}
void solve() {  // poly_exp (cppoly) is pasted above solve()
    int n, S;
    cin >> n >> S;
    vi cnt(S + 1);
    rep(i, n) {
        int w;
        cin >> w;
        cnt[w]++;
    }
    vi inv(S + 1, 1);  // inverses of 1..S modulo PMOD
    rep(i, 2, S + 1) inv[i] = (PMOD - PMOD / i) * inv[PMOD % i] % PMOD;
    Poly lg(S + 1, 0);  // log(1 + x^w) = x^w - x^2w / 2 + x^3w / 3 - ...
    rep(w, 1, S + 1) {
        if (!cnt[w]) continue;
        for (int j = 1; w * j <= S; j++) {  // the harmonic sum: O(S log S)
            ll t = cnt[w] % PMOD * inv[j] % PMOD;
            lg[w * j] = (lg[w * j] + (j % 2 ? t : PMOD - t)) % PMOD;
        }
    }
    Poly f = poly_exp(lg, S + 1);  // prod (1 + x^w) = exp(sum of the logs)
    print(vi(f.begin() + 1, f.end()));
}
\end{code}

\paragraph{Evaluating at many points.} $A(x_i)$ is the remainder of $A$ modulo $x - x_i$. Build the \emph{subproduct tree} whose leaves
are the factors $x - x_i$ and whose inner nodes multiply their children; then reduce $A$ modulo the root and push the remainders down,
each child taking the remainder modulo its own product (Figure~\ref{fig:subproduct}). Every level costs $\Oh{n \log n}$, so evaluating
at $n$ points costs $\Oh{n \log^2 n}$. Interpolation runs the same tree backwards, combining the Lagrange terms bottom up.

\begin{figure}[htbp]
\centering
\begin{forest}
  rectree, for tree={font=\sffamily\scriptsize, align=center, s sep=3mm, l sep=7mm, inner sep=3pt}
  [{$(x - x_1)(x - x_2)(x - x_3)(x - x_4)$\\[1pt]\textcolor{fgverm}{$A \bmod M$}}, hot
    [{$(x - x_1)(x - x_2)$\\[1pt]\textcolor{fgverm}{$r \bmod M_L$}}
      [{$x - x_1$\\[1pt]\textcolor{fgverm}{$A(x_1)$}}, good] [{$x - x_2$\\[1pt]\textcolor{fgverm}{$A(x_2)$}}, good]]
    [{$(x - x_3)(x - x_4)$\\[1pt]\textcolor{fgverm}{$r \bmod M_R$}}
      [{$x - x_3$\\[1pt]\textcolor{fgverm}{$A(x_3)$}}, good] [{$x - x_4$\\[1pt]\textcolor{fgverm}{$A(x_4)$}}, good]]]
\end{forest}
\caption{Multipoint evaluation. The products are built bottom up; the remainders (red) flow top down, each node reducing its parent's
remainder modulo its own product, until the leaves hold the values.}
\label{fig:subproduct}
\end{figure}

\begin{problem}{Multipoint evaluation}
Evaluate a polynomial of degree below $n$ at $m$ points modulo $998\,244\,353$ ($n, m \le 10^5$).
\solution \cc{poly_eval_many}. The constant is large: $10^5$ points take 1.4\,s, and below a few thousand points Horner's rule, $\Oh{nm}$,
is faster.
\end{problem}

\begin{code}
void solve() {  // poly_eval_many (cppoly) is pasted above solve()
    int n, m;
    cin >> n >> m;
    Poly a(n);
    vi xs(m);
    for (ll& c : a) cin >> c;
    for (ll& x : xs) cin >> x;
    print(poly_eval_many(a, xs));  // a subproduct tree: O((n + m) log^2)
}
\end{code}

\paragraph{Online convolution.} Some recurrences convolve a sequence with itself as it is being produced, such as $f_n = \sum_{i=1}^{n}
g_i f_{n-i}$ when $g_i$ itself depends on $f$, so the whole sequence is never available for one product. Divide and conquer on the
indices, exactly like CDQ (Section~\ref{sec:cdq}), solves this in $\Oh{n \log^2 n}$: compute $f_l, \dots, f_{m-1}$ recursively, add their
contributions to $f_m, \dots, f_{r-1}$ with one product of size $\Oh{r - l}$, then recurse on the right half (Figure~\ref{fig:onlineconv}).

\begin{figure}[htbp]
\centering
\begin{tikzpicture}[x=0.95cm]
  \foreach \ii in {0, ..., 7} \node[cell, minimum width=9.5mm] (F\ii) at (\ii, 0) {$f_{\ii}$};
  \draw[line width=0.8pt, draw=fgblue, rounded corners=2pt] (-0.05, -0.55) rectangle (3.95, -0.85);
  \draw[line width=0.8pt, draw=fgverm, rounded corners=2pt] (4.05, -0.55) rectangle (8.05, -0.85);
  \draw[line width=0.8pt, draw=fgblue!60, rounded corners=2pt] (-0.05, -1.1) rectangle (1.95, -1.4);
  \draw[line width=0.8pt, draw=fgverm!60, rounded corners=2pt] (2.05, -1.1) rectangle (3.95, -1.4);
  \draw[->, line width=1.1pt, draw=fgverm] (F2.north) to[bend left=40] node[above, font=\sffamily\scriptsize, text=fgverm]
    {$f_{0..3} * g$ adds to $f_{4..7}$} (F5.north);
  \node[notel, anchor=west, text width=4.6cm] at (8.6, -0.4) {First finish the left half (blue), recursively; then add its whole
    contribution to the right half (red) with one product; then recurse on the right half. Each level costs $\Oh{n \log n}$.};
\end{tikzpicture}
\caption{Online convolution by divide and conquer.}
\label{fig:onlineconv}
\end{figure}

% ======================================================================================================================
\section[Walsh--Hadamard transforms]{Walsh--Hadamard transforms \hfill \ioi}\label{sec:fwht}

Replace the sum of indices in a convolution by a bitwise operation: $c_k = \sum_{i \oplus j = k} a_i b_j$ is the \defn{XOR
convolution}, and $\land$ or $\lor$ give the AND and OR convolutions. Each has a transform that turns it into a pointwise product.
For XOR it is the \defn{Walsh--Hadamard transform}: process the bits one at a time, and for every pair of indices $(i, i + 2^b)$ that
differ only in bit $b$, replace $(u, v)$ by $(u + v, u - v)$. The transform of $a$ has the entry $\sum_i (-1)^{\mathrm{popcount}(i \land
k)} a_i$ at $k$, the sign is multiplicative over $\oplus$, and so the transform of the XOR convolution is the pointwise product. The
inverse is the same transform followed by a division by $n$ (Figure~\ref{fig:fwht}). For AND and OR the transforms are the sums over
supersets and over subsets of Section~\ref{sec:sosdp}, inverted by the M\"obius transform. Each costs $\Oh{n \log n}$ for $n = 2^k$.

\begin{figure}[htbp]
\centering
\begin{tikzpicture}
  \node[blbl, anchor=east] at (-0.2, 0) {$a$};
  \arrs{Fa}{0}{0}{1/n, 2/n, 0/n, 1/n}
  \node[blbl, anchor=east] at (-0.2, -0.8) {bit 0};
  \arrs{Fb}{0}{-0.8}{3/n, -1/n, 1/n, -1/n}
  \node[blbl, anchor=east] at (-0.2, -1.6) {bit 1: $\hat a$};
  \arrs{Fc}{0}{-1.6}{4/cool, -2/cool, 2/cool, 0/cool}
  \node[blbl, anchor=east] at (-0.2, -2.4) {$\hat a^2$};
  \arrs{Fd}{0}{-2.4}{16/n, 4/n, 4/n, 0/n}
  \node[blbl, anchor=east] at (-0.2, -3.2) {inverse: $a \oplus a$};
  \arrs{Fe}{0}{-3.2}{6/good, 4/good, 4/good, 2/good}
  \node[notel, anchor=west, text width=8.6cm] at (3.4, -1.6) {The values 0, 1, 1, 3 as a histogram $a$ over $\{0, 1, 2, 3\}$. Each
    pass replaces the pairs that differ in one bit by $(u + v, u - v)$. Squaring the transform and transforming back (then dividing
    by 4) gives the XOR convolution of $a$ with itself: among the 16 ordered pairs, 6 have XOR 0, 4 have XOR 1, 4 have XOR 2 and 2 have
    XOR 3.};
\end{tikzpicture}
\caption{The Walsh--Hadamard transform of a 4-entry histogram, and the XOR convolution it computes.}
\label{fig:fwht}
\end{figure}

\usetemplate{Walsh--Hadamard transform}{Ctrl+Alt+U, W}{cpfwht}

\begin{problem}{XOR pairs}
For $n \le 2 \cdot 10^5$ values below $2^k$ ($k \le 17$) and every $x < 2^k$, count the ordered pairs $(i, j)$ with $a_i \oplus a_j =
x$, modulo $10^9 + 7$.
\solution The XOR convolution of the histogram with itself: $\Oh{2^k k}$, 0.10\,s.
\end{problem}

\begin{code}
void solve() {  // xor_convolution (cpfwht) is pasted above solve()
    int n, k;
    cin >> n >> k;
    vi cnt(1 << k);
    rep(i, n) {
        int x;
        cin >> x;
        cnt[x]++;
    }
    print(xor_convolution(cnt, cnt));  // ordered pairs with a_i ^ a_j = x
}
\end{code}

% ======================================================================================================================
\section[Linear recurrences]{Linear recurrences \hfill \ioi}\label{sec:linrec}

A sequence satisfies a \defn{linear recurrence} of order $L$ if $a_n = c_1 a_{n-1} + c_2 a_{n-2} + \dots + c_L a_{n-L}$ for all $n \ge
L$. Section~\ref{sec:matexp} computed $a_n$ with an $L \times L$ matrix power in $\Oh{L^3 \log n}$. Two better tools exist.

\paragraph{Kitamasa's method.} Shifting the sequence by one step is multiplication by $x$ modulo the characteristic polynomial $P(x) =
x^L - c_1 x^{L-1} - \dots - c_L$, because $P$ encodes exactly the rule that rewrites $a_L$ in terms of earlier terms. So compute $x^n
\bmod P(x) = \sum_{i<L} r_i x^i$ by repeated squaring of polynomials modulo $P$; then $a_n = \sum_i r_i a_i$. Each squaring is a product
of two polynomials of degree below $L$ and a reduction, $\Oh{L^2}$ directly or $\Oh{L \log L}$ with the polynomial toolkit, so $a_n$
costs $\Oh{L^2 \log n}$ instead of $\Oh{L^3 \log n}$.

\paragraph{Berlekamp--Massey.} Conversely, given the first terms of a sequence, the \textbf{Berlekamp--Massey algorithm} finds the
shortest linear recurrence that generates them, in $\Oh{N^2}$ for $N$ terms. It processes the terms one by one, keeping the current
recurrence; when the recurrence predicts a term wrongly, it corrects itself by adding a shifted multiple of the recurrence it had at its
last failure, which cancels the error without breaking the earlier terms. If the sequence truly satisfies a recurrence of order $L$,
$2L$ terms are enough to find it.

\usetemplate{Berlekamp--Massey / Kitamasa}{Ctrl+Alt+U, Shift+B / Ctrl+Alt+U, K}{cpbm, cpkitamasa}

\begin{trick}[Guess the recurrence]
Many counting sequences satisfy a linear recurrence whose order is small but tedious to derive: DPs whose state is a small vector (tilings
of a $k \times n$ board, walks in a fixed graph, automata), and sums of such sequences. Compute the first few hundred terms with a simple
DP, run Berlekamp--Massey modulo a prime, and use Kitamasa for $n = 10^{18}$. If the recurrence found has order close to half the number
of terms, it is probably not the true one: compute more terms and check.
\end{trick}

\begin{problem}{Throwing Dice, again (CSES)}
Count the ways to reach a sum of $n \le 10^{18}$ with dice throws, modulo $10^9 + 7$.
\solution Section~\ref{sec:matexp} derived the recurrence of order 6 and used a matrix power. Without any derivation, 20 terms from the
obvious DP let Berlekamp--Massey find $f_n = f_{n-1} + \dots + f_{n-6}$ by itself, and Kitamasa jumps to $n$: 0.06\,s.
\end{problem}

\begin{code}
void solve() {  // berlekamp_massey (cpbm), linear_rec_nth (cpkitamasa) pasted
    ll n;
    cin >> n;
    vi f(20);  // the first terms by the DP f[s] = f[s - 1] + ... + f[s - 6]
    f[0] = 1;
    rep(s, 1, 20) rep(d, 1, 7) if (s >= d) f[s] = (f[s] + f[s - d]) % MOD;
    vi c = berlekamp_massey(f);  // recovers the recurrence from the terms alone
    print(linear_rec_nth(f, c, n));
}
\end{code}

\begin{summary}[Chapter summary]
\begin{itemize}
  \item $A^k$ counts walks with $k$ edges; over $(\min, +)$ it finds the cheapest such walks, over $(\lor, \land)$ reachability.
  \item Gauss--Jordan elimination solves linear systems modulo a prime in $\Oh{nm \min(n, m)}$ and reveals the rank, the free variables
        and the number of solutions; over the reals, pivot on the largest entry.
  \item A linear basis over GF(2) answers maximum, counting and membership questions about subset XORs; prefix bases that prefer recent
        vectors answer them for ranges.
  \item The FFT multiplies polynomials in $\Oh{n \log n}$ through evaluation at roots of unity; the NTT does it exactly modulo
        $998\,244\,353$; reversing a sequence turns convolution into correlation.
  \item Newton's method doubles the precision of power series inverses, logarithms, exponentials and roots, each $\Oh{n \log n}$; a
        subproduct tree evaluates and interpolates at many points in $\Oh{n \log^2 n}$.
  \item Walsh--Hadamard and subset-sum transforms turn XOR, AND and OR convolutions into pointwise products.
  \item Kitamasa computes the $n$-th term of a recurrence in $\Oh{L^2 \log n}$, and Berlekamp--Massey finds the recurrence from $2L$
        terms.
\end{itemize}
\end{summary}

\begin{practice}
\ppgroup{Linear algebra and transforms}
\begin{pplist}
  \item Fibonacci Numbers
  \item Throwing Dice
  \item Graph Paths I
  \item Graph Paths II
  \item System of Linear Equations
  \item Apples and Bananas
  \item One Bit Positions
  \item Signal Processing
\end{pplist}
\end{practice}
